\subsection{Definizione di automa a stati finiti.}

\begin{defi}\label{def:1}
Siano $A$ e $Q$ insiemi finiti, un \textit{automa finito} di alfabeto $A$ e insieme degli stati $Q$  è una tripla $(T,q_0,F)$ dove
\[
T: Q\times A \to Q
\]
$F \subset Q$ e $q_0\in Q$ è uno stato particolare, detto stato iniziale.
\end{defi}


\subsection{Definizione di algoritmo associato ad un automa}

A partire dalla tripla $(T,q_0,F)$ che specifica un dato automa di alfabeto $A$ e insieme degli stati $Q$ si costruisce prima l'insieme delle configurazioni possibili:

\begin{defi}[Configurazione per Automa Finito]\label{def:3}
Una configurazione per automa a stati finiti di alfabeto $A$ e insieme degli stati $Q$ è una tripla $(w,p,q)$, dove
\begin{enumerate}
    \item $w\in A^*$ è la parola da analizzare,
    \item $p\in\mathbb{N}$, è un intero che indica la posizione della parola da analizzare nella prossima transizione,
    \item $q\in Q$ indica lo stato corrente e costituisce la memoria delle transizioni effettuate nei precedenti passi di computazione.
\end{enumerate}
\end{defi}

L'insieme delle configurazioni di alfabeto $A$ e insieme degli stati $Q$ per automa finito sarà denotato con
\[
Conf_{DFA}^{A,Q} = A^* \times \mathbb{N} \times Q
\]
ed allora possiamo definire:

\begin{defi}[Algoritmo associato ad un automa finito deterministico]\label{def:4}
L'algoritmo associato ad un automa di alfabeto $A$ e insieme degli stati $Q$ è l'algoritmo $$(A,Conf_{DFA}^{A,Q}, \{0,1\},\epsilon,\tau,\delta)$$ dove le tre funzioni sono definite come segue:
\begin{itemize}
    \item la \textbf{funzione di ingresso} che associa ad una parola $w\in A^*$, la configurazione $\varepsilon(w) = (w,1,q_0)$,

    \item la \textbf{funzione di transizione} $\tau: Conf_{DFA}^{A,Q} \to Conf_{DFA}^{A,Q}$ è data da
    \[
    \tau((w,p,q)) = (w,p+1,T(q,w(p))),
    \]

    \item la \textbf{funzione di uscita} $\delta: Conf_{DFA}^{A,Q} \to \{0,1\}*$, data da
    \[
    \delta((w,p,q)) =
    \begin{cases}
    0 & \text{se } p>|w| \text{ e se } q\notin F, \\
    1 & \text{se } p>|w| \text{ e se } q\in F.
    \end{cases}
    \]
\end{itemize}
\end{defi}



\begin{defi}[Insiemi decidibili per automa]\label{def:5}
Un insieme $X\subset A^*$, è \textit{decidibile} per automa finito (oppure \textit{DFA-decidibile} oppure \textit{regolare}), se esiste un insieme finito $Q$ e un automa finito con alfabeto $A$ e insieme degli stati $Q$ tale che l'algoritmo associato decide $X$.
\end{defi}
\begin{exam}\label{ex: parole lunghezze pari}
Descrivere l'automa finito che decide l'insieme delle parole di lunghezza pari (risp. dispari) su un alfabeto dato $A$.

A partire dall'alfabeto finito $A$, e da un insieme di due stati $Q=\{0,1\}$ che corrisponderanno rispettivamente allo stato ``abbiamo letto un numero pari (oppure dispari) di caratteri della parola''; costruiamo la funzione di transizione che ad ogni carattere letto cambia stato:
\[
T: Q\times A \to Q
\]
tale che
\[
T(q,a) =
\begin{cases}
1 & \text{se } q=0 \\
0 & \text{se } q=1,
\end{cases}
\]
per ogni $a\in A$.

Avremo allora che lo stato iniziale $q_0=0$, poiché inizialmente non avremo ancora letto nessun carattere e che l'insieme degli stati accettanti sarà: $F=\{0\}$.
\begin{proof}
Sia
\[
X=\{\,w\in A^* : |w|\text{ è pari}\,\}.
\]
Dimostriamo che l'algoritmo associato all'automa
\[
(T,q_0,F)
\]
decide l'insieme \(X\).

Sia dunque \(w=a_1a_2\cdots a_n\in A^*\). La configurazione iniziale è
\[
\epsilon(w)=(w,1,0).
\]

Ad ogni applicazione della funzione di transizione viene letta una lettera della parola e lo stato viene invertito:
\[
0\longleftrightarrow1.
\]
Per induzione sul numero di transizioni effettuate si ottiene che dopo aver letto i primi \(k\) caratteri della parola la configurazione è
\[
(w,k+1,q_k),
\]
dove
\[
q_k=
\begin{cases}
0 & \text{se } k \text{ è pari},\\
1 & \text{se } k \text{ è dispari}.
\end{cases}
\]

Dopo la lettura dell'intera parola si raggiunge quindi la configurazione
\[
(w,n+1,q_n).
\]
Poiché
\[
n+1>|w|,
\]
questa appartiene al dominio della funzione di uscita. Inoltre
\[
q_n=
\begin{cases}
0 & \text{se } |w| \text{ è pari},\\
1 & \text{se } |w| \text{ è dispari}.
\end{cases}
\]

Essendo \(F=\{0\}\), segue che
\[
\delta(w,n+1,q_n)=
\begin{cases}
1 & \text{se } |w| \text{ è pari},\\
0 & \text{se } |w| \text{ è dispari}.
\end{cases}
\]

Pertanto
\[
f_{\mathcal A}(w)=
\begin{cases}
1 & \text{se } w\in X,\\
0 & \text{se } w\notin X,
\end{cases}
\]
cioè
\[
f_{\mathcal A}=\mathbf 1_X.
\]

Ne consegue che l'algoritmo associato all'automa decide l'insieme delle parole di lunghezza pari.
\end{proof}
\end{exam}
\begin{ex}\label{es:1}
Fissiamo l'alfabeto binario $A=\{0,1\}$, dimostrare che l'insieme
\[
X = \{ w0 \mid w\in A^* \}
\]
è decidibile per automa a stati finiti.


\textbf{Soluzione.} Si richiede di mostrare un automa il cui algoritmo associato accetta le parole che terminano con uno zero. È dunque sufficiente memorizzare nello stato l'ultimo simbolo letto e alla fine della lettura della parola decidere in base allo stato: poniamo $Q=\{q_0,q_1\}$, e definiamo
\[
T(q,a) :=
\begin{cases}
q_0 & \text{se } a=0, \\
q_1 & \text{se } a=1,
\end{cases}
\]
prendiamo $q_0$ come stato iniziale e $F=\{q_0\}$ come insieme degli stati finali accettanti.
\end{ex}
\begin{proof}
Sia
\[
X=\{w0\mid w\in A^*\}
\]
l'insieme delle parole sull'alfabeto binario che terminano con il simbolo \(0\).

Costruiamo un automa finito che decide \(X\). Consideriamo l'insieme degli stati
\[
Q=\{q_0,q_1\},
\]
dove intuitivamente:
\begin{itemize}
    \item \(q_0\) rappresenta la situazione in cui l'ultimo simbolo letto è \(0\);
    \item \(q_1\) rappresenta la situazione in cui l'ultimo simbolo letto è \(1\).
\end{itemize}

Definiamo la funzione di transizione
\[
T:Q\times A\to Q
\]
come
\[
T(q,a)=
\begin{cases}
q_0 & \text{se } a=0,\\
q_1 & \text{se } a=1.
\end{cases}
\]

Scegliamo come stato iniziale
\[
q_{\mathrm{init}}=q_0
\]
e come insieme degli stati accettanti
\[
F=\{q_0\}.
\]

Consideriamo ora l'algoritmo associato all'automa.
Sia
\[
w=a_1a_2\cdots a_n\in A^*.
\]
La configurazione iniziale è
\[
\epsilon(w)=(w,1,q_0).
\]

Dopo \(k\) transizioni, per induzione su \(k\), la configurazione raggiunta è
\[
(w,k+1,q_k),
\]
dove
\[
q_k=
\begin{cases}
q_0 & \text{se } a_k=0,\\
q_1 & \text{se } a_k=1.
\end{cases}
\]

Infatti, ad ogni passo la funzione di transizione sostituisce lo stato corrente con lo stato corrispondente all'ultimo carattere letto.

Dopo aver letto tutta la parola \(w\), la computazione termina nella configurazione
\[
(w,n+1,q_n).
\]
La funzione di uscita restituisce quindi
\[
\delta(w,n+1,q_n)=
\begin{cases}
1 & \text{se } q_n\in F,\\
0 & \text{se } q_n\notin F.
\end{cases}
\]

Poiché
\[
F=\{q_0\},
\]
si ha
\[
\delta(w,n+1,q_n)=
\begin{cases}
1 & \text{se l'ultimo simbolo di } w \text{ è }0,\\
0 & \text{se l'ultimo simbolo di } w \text{ è }1.
\end{cases}
\]

Quindi
\[
f_{\mathcal A}(w)=
\begin{cases}
1 & \text{se } w\in X,\\
0 & \text{se } w\notin X.
\end{cases}
\]

Pertanto
\[
f_{\mathcal A}=\mathbf 1_X
\]
e quindi l'algoritmo associato all'automa decide \(X\).

Concludiamo che
\[
X=\{w0\mid w\in A^*\}
\]
è DFA-decidibile.
\end{proof}